\chapter{Phá Băng Bật Nói}
\chapterintro{Câu hỏi đầu tiên quyết định nhịp lớp}{Nếu mở đúng, học sinh sẽ nói tự nhiên hơn cả tiết.}

\section{Nếu tiết học này là một món ăn, nó sẽ là món gì?}
\begin{cauhoibox}{Câu hỏi}
Nếu tiết học này là một món ăn, nó sẽ là món gì?
\end{cauhoibox}
\begin{taisaocoich}
Câu hỏi này không có đúng sai nên học sinh đỡ sợ. Nó vừa vui, vừa mở được liên tưởng, lại đủ nhẹ để cả lớp nói nhanh.
\end{taisaocoich}
\begin{goiybien}
Đổi thành: màu gì, bài hát gì, thời tiết gì. Hợp với đầu tiết hoặc lúc cần kéo lớp tỉnh lên.
\end{goiybien}
\ngancach

\section{Ai có một câu trả lời kỳ cục nhưng hợp lý?}
\begin{cauhoibox}{Câu hỏi}
Ai có một câu trả lời kỳ cục nhưng hợp lý?
\end{cauhoibox}
\begin{taisaocoich}
Câu này cho phép học sinh sáng tạo mà không sợ bị cười vì nói khác. Nó rất hợp để bật lớp trong 1 phút đầu.
\end{taisaocoich}
\begin{goiybien}
Nhớ thêm luật nhỏ: kỳ cục nhưng phải giải thích được. Thế là lớp vừa vui vừa không bị trôi hẳn.
\end{goiybien}
\ngancach

\section{Nếu em được đặt tên mới cho bài hôm nay, em đặt là gì?}
\begin{cauhoibox}{Câu hỏi}
Nếu em được đặt tên mới cho bài hôm nay, em đặt là gì?
\end{cauhoibox}
\begin{taisaocoich}
Câu hỏi này khiến học sinh nhìn bài theo góc của mình. Rất tốt để chuyển từ trạng thái nghe thụ động sang tham gia chủ động.
\end{taisaocoich}
\begin{goiybien}
Cho 3 em đặt tên, rồi dùng tên vui nhất làm tên phụ của tiết học trên bảng.
\end{goiybien}
\ngancach

\section{Nếu lớp mình có một nút mở lời, em bấm chỗ nào?}
\begin{cauhoibox}{Câu hỏi}
Nếu lớp mình có một nút mở lời, em bấm chỗ nào để ai cũng muốn nói hơn?
\end{cauhoibox}
\begin{taisaocoich}
Câu này khiến học sinh nghĩ về chính không khí lớp học. Nó vừa là phá băng, vừa cho thầy cô biết lớp thích kiểu mở đầu nào.
\end{taisaocoich}
\begin{goiybien}
Có thể đổi thành: nút tỉnh lại, nút bớt ngại, hoặc nút bớt im. Rất hợp đầu giờ hoặc đầu buổi sinh hoạt.
\end{goiybien}
\ngancach

\section{Một câu trả lời nào sẽ làm bạn bên cạnh muốn nói tiếp?}
\begin{cauhoibox}{Câu hỏi}
Một câu trả lời nào sẽ làm bạn bên cạnh muốn nói tiếp?
\end{cauhoibox}
\begin{taisaocoich}
Thay vì chỉ hỏi để một em trả lời, câu này kéo học sinh nghĩ tới việc mở lời cho nhau. Nó rất tốt để bật lớp theo kiểu lan dần.
\end{taisaocoich}
\begin{goiybien}
Sau một câu trả lời đầu tiên, mời ngay bạn bên cạnh nối tiếp để lớp thấy câu hỏi đang thật sự mở ra thêm tiếng nói.
\end{goiybien}
